% Generated draft: requires mathematical and source review.
\documentclass[11pt]{article}
\usepackage[margin=25mm]{geometry}
\usepackage{fontspec}
\setmainfont{Latin Modern Roman}
\usepackage{amsmath,amssymb,mathtools,unicode-math}
\setmathfont{Latin Modern Math}
\usepackage{enumitem,needspace,xurl,hyperref}
\setlength{\parindent}{0pt}
\setlength{\parskip}{5pt}
\setlength{\emergencystretch}{3em}
\newcommand{\N}{\mathbb{N}}
\newcommand{\Z}{\mathbb{Z}}
\newcommand{\Q}{\mathbb{Q}}
\newcommand{\R}{\mathbb{R}}
\newcommand{\C}{\mathbb{C}}
\newcommand{\F}{\mathbb{F}}
\newcommand{\A}{\mathbb{A}}
\newcommand{\PP}{\mathbb{P}}
\newcommand{\E}{\mathbb{E}}
\newcommand{\eps}{\varepsilon}
\newcommand{\dd}{\,\mathrm{d}}
\begin{document}
\section*{UnsolvedMath ID: 30002718}
\textbf{Source problem code:} catalogue entry
\subsection{Definitions and mathematical statement}
Let \(\mathbb H^2\) denote the hyperbolic plane, with hyperbolic distance \(d_{\mathbb H}\) and area measure \(A_{\mathbb H}\). The supplied record does not specify the curvature normalization. A point configuration \(X\subseteq\mathbb H^2\) is locally finite if every compact subset contains finitely many points of \(X\). A Poisson point process with locally finite intensity measure \(\mu\) is a random locally finite configuration whose counts in pairwise disjoint bounded Borel sets are independent, with
\[
\#(X\cap B)\sim\operatorname{Poisson}(\mu(B)).
\]
The homogeneous choice \(\mu=\lambda A_{\mathbb H}\), \(\lambda>0\), is one possible specification; the record does not explicitly select it. In a finite-point model, \(X\) is a random finite configuration, whose sampling region, point-count law, and joint location law must be specified.

For a specified random graph construction \(G=(X,E)\), let \(C_G(x)\) be the connected component containing \(x\in X\), and let
\[
N_k(G)=\#\{C:C\text{ is a connected component of }G,\ |C|=k\},\qquad k\geq1.
\]
For finite nonempty \(X\), define
\[
L(G)=\max_{x\in X}|C_G(x)|.
\]
Then \(G\) is connected precisely when \(L(G)=|X|\). For an infinite locally finite configuration, the graph percolation event is
\[
\mathcal P_G=\{\exists x\in X:\ |C_G(x)|=\infty\}.
\]
If continuum percolation instead refers to an occupied random subset of \(\mathbb H^2\), its construction must be given, and the corresponding event is the existence of an unbounded connected component of that subset. The record does not specify which continuum construction is intended.

\textbf{Research question, conditional on specifying the models.} Determine the component-size and connectivity behavior of finite random geometric graphs and continuum percolation in \(\mathbb H^2\), when hyperbolic geometry replaces Euclidean planar geometry. For a chosen graph model, possible precise observables include the laws of \(N_k(G)\), \(L(G)\), and the connectivity event, and, for an infinite model, the probability of \(\mathcal P_G\). These observables provide a framework for the stated question, rather than additional source-specified conjectures.

\textbf{Unresolved specification.} The source supplies no connection rule, connection range or connection function, intensity or finite sampling law, observation region, boundary convention, parameter scaling, limiting regime, or proposed quantitative conclusion. It also does not identify the exact Euclidean statements to be transferred. Consequently, the supplied question is a research program and does not determine a unique quantified theorem or conjecture. The definitions above do not select any of these missing model assumptions.
\subsection{Short English statement}
The question asks how component sizes and connectivity change when random geometric graphs and continuum percolation are studied in the hyperbolic plane. A precise conjecture requires further choices of point distribution, connection rule, boundary conditions, and limiting regime. The supplied partially solved label is a dated curation assessment, not an independently verified mathematical conclusion.
\subsection{Sources}
\begin{itemize}
\item Local UnsolvedMath registry, \texttt{data/unsolvedmath/problems.json}, permanent ID \(30002718\), title ``Random Geometric Graphs in the Hyperbolic Plane,'' exact source code \texttt{catalogue entry}; supplied complete source record from the Oberwolfach Reports Open Problems collection. No individual proposer is supplied; the proposed year is recorded as 2014.
\item Supplied report citation: \emph{Probability, Trees and Algorithms (2015)}, Oberwolfach Reports, pp.\ 2825--2871, DOI \texttt{10.4171/owr/2014/50}. \url{https://doi.org/10.4171/owr/2014/50}. This is also the supplied literature-source entry labeled ``Oberwolfach report.'' The report title's year 2015 and the record's proposed year 2014 are retained without resolving their relationship.
\item Supplied literature source: ``Benjamini--Schramm, hyperbolic percolation.'' \url{https://arxiv.org/abs/math/9912233}.
\item Supplied literature source: ``Hyperbolic random connection model.'' \url{https://doi.org/10.1112/jlms.70345}.
\item Literature assessment checked at \texttt{2026-08-21}: the supplied curation labels the problem partially solved and reports rigorous phase diagrams and recent critical-exponent results for hyperbolic percolation and random-connection models, while describing the broad question across the recorded variants as not fully settled. This is supplied dated curation, not an independently verified literature assessment; time-sensitive status and cited claims require independent verification.
\item The registry labels the statement \texttt{corrected\_verified} and records a manual semantic, typographic, and LaTeX review asserting a meaning-preserving, self-contained cleanup. Those are supplied editorial claims and do not establish a precise mathematical formulation or mathematical correctness.
\end{itemize}
\end{document}
